\chapter*{Extended Abstract}
In this thesis, a preliminary hexagonal geometry of the  SMR model of the ALFRED (Advanced Lead Fast REactor Demonstrator) concept design for a lead cooled reactor is analyzed with a computational platform that integrates different computational tools into the common framework given by the SALOME platform software. The purpose of the paper is to investigate the multiphysics (thermal-hydraulic-neutronic) approach to this SMR model by coupling the DRAGON-DONJON codes, for lattice calculations and full core simulation, with a 3D-porous media thermal-hydraulic code, FEMUS. 

In particular, the lattice code, DRAGON, is used to evaluate the macroscopic cross sections over the hexagonal lattice devised for ALFRED, collapsing the microscopic cross section data into 33 groups, parametrized on fuel and moderator temperature, density and power distribution. With the macroscopic cross sections for the various lattice cells of the reactor defined, the distribution of neutron flux in the core is obtained by the full core simulation with the DONJON code, that may be used in the simulation of fuel management, reactor operation or accident scenarios. 

In transient simulations (in a quasi-static approach), at each time step, the neutron flux is computed with the corresponding thermal source while the thermal-hydraulic module (FEMUS) computes the distribution of the coolant velocity, pressure and temperature fields in the reactor core. Then the neutron modules can receive all the feedback variables, as fuel temperature, moderator temperature and density that can be cast in input to the DRAGON code, where macroscopic cross sections are interpolated and total neutron fluxes evaluated. 

In this way, one can obtain an iterative process for studying the transient evolution of the model of the ALFRED fast lead cooled reactor as SMR case-study and preliminary results about his time-dependent behavior can be fully analyzed.